% ===================================================================
%  तक्ता क्रमांक १३ — हॉटेल. --- उपाहारगृह. कॅफे.
%
%  Traduction, faite en 2026, en marathi d'aujourd'hui, sur l'ido de
%  texto/io. Regles de la colonne : voir 10-takta-01.tex. Un seul
%  recit a la premiere personne, alineas numerotes : les cles ne
%  portent pas d'indice de scene. Les DEUX notes d'auteur sont en
%  gras-italique dans le fac-simile, la ou celles des tableaux 1, 5
%  et 8 sont en romain : la traduction garde la difference.
%
%  LES ECHECS REVIENNENT CHEZ EUX, ET C'EST LE MOT LE PLUS ANCIEN DE
%  LA COLONNE. L'ido ecrit « shako », emprunte au persan شاه par
%  l'italien et le francais ; le marathi dit बुद्धिबळ — la force de
%  l'esprit — et derriere lui il y a चतुरंग, le jeu des quatre
%  membres de l'armee, d'ou tout le reste est parti. Le livret
%  emprunte a l'Europe un mot que l'Europe avait emprunte a l'Inde.
%  Aucun autre jeu de la liste n'est dans ce cas : le trictrac et
%  les dames s'ecrivent बॅकगॅमन et चेकर्स, parce que ce ne sont pas
%  les memes jeux que सोंगट्या ou चौसर et qu'on ne transpose pas les
%  choses.
%
%  L'HOTEL, LUI, EST TENU PAR DES MOTS MARATHES : गल्ला la caisse,
%  पोऱ्या le chasseur, द्वारपाल le portier, आचारी le cuisinier,
%  परटीण la blanchisseuse — du nom de la caste des परीट —, हमाल le
%  porteur, वाढपी le garcon, बक्षिशी le pourboire, वाचनालय le salon
%  de lecture, लिफाफा l'enveloppe, शाई l'encre, पत्ते les cartes,
%  चाबूक le fouet, पताका le pavillon.
%
%  ET LE MENU SUIT LA REGLE DES CHOSES, PAS CELLE DE LA LANGUE. Les
%  crevettes sont कोळंबी et les epinards पालक — le Maharashtra les
%  connait —, mais les tripes a la mode de Caen restent les tripes a
%  la mode de Caen. Traduire le repas en repas d'ici donnerait une
%  phrase plus appetissante et une traduction fausse.
%
%  DEUX INVERSIONS PREVENUES, ET NEUF QUI ONT PASSE — le plus mauvais
%  compte de la colonne. La cause est visible : ce tableau est une
%  seule longue enumeration de rapports (« la chambre du deuxieme
%  etage », « le commis de l'hotel », « la corbeille de la
%  blanchisseuse », « la caisse de la boutique »), et la
%  postposition mord sur chacun. La ou le tableau 11 alignait des
%  batiments cote a cote et ne coutait que trois inversions, celui-ci
%  ne fait qu'accrocher des choses les unes aux autres.
%
%  ET L'UNE DES DEUX « PREVENUES » N'ETAIT QU'A MOITIE FAITE : en
%  redressant le conducteur de l'omnibus (30, 29) j'avais laisse
%  passer la malle du voyageur (31, 32), qui est dans la meme phrase
%  et dans le meme piege. Corriger un rapport ne corrige pas son
%  voisin.
% ===================================================================

%%K t13-noto-1 noto
\VUnotes{143.2mm}{%
\textit{\VUgras{(*) खोलीच्या तपशिलांसाठी तक्ता क्रमांक ६ पाहा.}}%
}

%%K t13-apar-1 apar
\VUsaut{3.0mm}
\VUtitre{15.0pt}{60}{तक्ता क्रमांक १३}
\VUsaut{1.6mm}
\VUpk{10.2pt}{40}{हॉटेल. --- उपाहारगृह.}
\VUsaut{2.0mm}
\VUpk{10.2pt}{40}{कॅफे.}
\VUsaut{3.4mm}

%%K t13-01-1 p
१. --- काल संध्याकाळी रुआयाँला पोहोचल्यावर मी « \textit{महासागर हॉटेलात} » उतरलो, जे एका
मित्राने मला सुचवले होते.

%%K t13-01-2 p
मला एक सुंदर \VUgras{खोली} \textsuperscript{(2)} दिली, दुसऱ्या \VUgras{मजल्यावर}
\textsuperscript{(1)}, जिथे मी फार छान झोपलो (*).

%%K t13-02-1 p
२. --- आज सकाळी, ज्या खोलीत \VUgras{नोकराणीने} \textsuperscript{(3)} न्याहारी आणून दिली होती
तिच्यातून बाहेर पडून, मी \VUgras{धूम्रपानाच्या खोलीत} \textsuperscript{(5)} शिरलो, जी
\VUgras{वाचनालय} \textsuperscript{(4)} म्हणूनही वापरतात — \VUgras{वर्तमानपत्रे}
\textsuperscript{(6)} वाचायला, \VUgras{सिगारेट} \textsuperscript{(8)} ओढत. वाचून झाल्यावर मी
\VUgras{उद्वाहकाने} \textsuperscript{(9)} खाली आलो आणि शहरात फिरायला गेलो.

%%K t13-03-1 p
३. --- हॉटेलच्या \VUgras{कारकुनाला} \textsuperscript{(12)} विचारायचे मी विसरलो होतो की
\VUgras{सामाईक जेवणाच्या मेजावर} \textsuperscript{(11)} जेवण कोणत्या वेळी वाढतात, म्हणून मी
लवकर परतलो आणि त्याचा फायदा घेऊन हॉटेलच्या \VUgras{न्हाणीघरात} \textsuperscript{(10)} आंघोळ
करायला गेलो. मग मी माझ्या एका मित्राला दूरध्वनी करायला पुन्हा \VUgras{गल्ल्याकडे}
\textsuperscript{(15)} गेलो. पण \VUgras{दूरध्वनी} \textsuperscript{(13)} वापरात होता; माझी
अडचण पाहून \VUgras{गल्लेवालीने} \textsuperscript{(14)}, जी एक \VUgras{बीजक}
\textsuperscript{(17)} \VUgras{प्रवासीनोंदवहीवर} \textsuperscript{(16)} नोंदवत होती,
\VUgras{द्वारपालाला} \textsuperscript{(18)} विनंती केली की \VUgras{पोऱ्याला}
\textsuperscript{(19)} पाठवून माझा निरोप \VUgras{सायकलवरून} \textsuperscript{(20)} पोहोचवावा.

%%K t13-04-1 p
४. --- मी तिचे आभार मानले आणि त्या \VUgras{हमालाला} \textsuperscript{(24)} (मजुराला) बक्षिशी
द्यायला बाहेर पडलो, ज्याने स्थानकाहून आपल्या \VUgras{हातगाडीवरून} \textsuperscript{(25)}
\VUgras{माझी टोपीची पेटी} \textsuperscript{(26)}, \VUgras{माझी सुटकेस} \textsuperscript{(27)}
आणि \VUgras{माझी प्रवासाची पिशवी} \textsuperscript{(28)} आणली होती. नुकत्याच आलेल्या
\VUgras{पत्रे वाटणाऱ्याने} \textsuperscript{(21)} मला एक \VUgras{पत्र} \textsuperscript{(22)}
दिले.

%%K t13-05-1 p
५. --- रस्त्यावरची वर्दळ पाहायला मी दाराच्या उंबरठ्यावर \VUgras{पत्रपेटीजवळ}
\textsuperscript{(23)} आणखी थोडा वेळ थांबलो. \VUgras{वाहक} \textsuperscript{(30)},
\VUgras{घोडागाडीचा} \textsuperscript{(29)}, आपल्या गाडीवर एक \VUgras{ट्रंक}
\textsuperscript{(31)} चढवत होता — एका \VUgras{प्रवाशाची} \textsuperscript{(32)}, ज्याला
\VUgras{हॉटेलमालक} \textsuperscript{(33)} निरोप देत होता. एका तरुण \VUgras{तारवाल्याने}
\textsuperscript{(35)} निवांतपणे हॉटेलच्या एका गिऱ्हाइकासाठी \VUgras{तार}
\textsuperscript{(34)} आणली; एक \VUgras{पर्यटक} \textsuperscript{(36)} खांद्याला
\VUgras{छायाचित्रयंत्र} \textsuperscript{(38)} अडकवून आपली \VUgras{मार्गदर्शिका}
\textsuperscript{(37)} पाहत होता.

%%K t13-tit-1 sub
\VUsaut{4.0mm}
\VUpk{10.2pt}{40}{न्याहारीची वेळ.}
\VUsaut{3.0mm}

%%K t13-06-1 p
६. --- जाळीसमोर एक निश्चिंत \VUgras{निरोप्या} \textsuperscript{(39)} आपला
\VUgras{ओझ्याचा आकडा} \textsuperscript{(40)} आणि आपली \VUgras{बूटपॉलिशची पेटी}
\textsuperscript{(41)} यांच्यामध्ये डुलकी घेत होता, तर एक तरुण \VUgras{परटीण}
\textsuperscript{(42)}, जिने धुण्याची \VUgras{टोपली} \textsuperscript{(43)} घेतली होती, ती
दाराशी उभ्या असलेल्या \VUgras{वाढप्यांच्या प्रमुखाच्या} \textsuperscript{(44)} गंभीर
चेहऱ्याकडे पाहून हसत होती.

%%K t13-07-1 p
७. --- \VUgras{आचाऱ्याला} \textsuperscript{(45)} पाहून — जो \VUgras{स्वयंपाकघरातून}
\textsuperscript{(47)} बाहेर आला होता, \VUgras{तळघरातल्या मजल्यावरच्या}
\textsuperscript{(48)}, \VUgras{भांडी घासणाऱ्याला} \textsuperscript{(46)} निरोप द्यायला
पाठवण्यासाठी — मला आठवले की न्याहारीची वेळ जवळ आली, आणि मी अंगण ओलांडून उपाहारगृहात शिरलो;
त्या अंगणात एक \VUgras{मोटारचालक} \textsuperscript{(50)} आपली \VUgras{मोटार}
\textsuperscript{(49)} उभी करत होता, तर एक \VUgras{यंत्रकारागीर} \textsuperscript{(52)}
\VUgras{सायकलस्वाराच्या} \textsuperscript{(53)} सूचनांप्रमाणे नुकतीच अपघात झालेली
\VUgras{पेट्रोलची तीनचाकी} \textsuperscript{(51)} दुरुस्त करत होता.

%%K t13-08-1 p
८. --- एका तरुण \VUgras{पोऱ्याला} \textsuperscript{(54)}, जो \VUgras{बिअरचे पेले}
\textsuperscript{(55)} घेऊन जात होता, त्याला मी विचारले की सामाईक मेज व्हरांड्याखाली आहे का,
आणि त्याच्या होकारार्थी उत्तरावरून मी एका मेजाशी जागा घेतली आणि उत्तम जेवण मागवले.

%%K t13-noto-2 noto
\VUnotes{143.2mm}{%
\textit{\VUgras{(*) शब्दकोशात दिलेल्या पदार्थांच्या यादीच्या मदतीने शिक्षकाला खालची यादी सहज बदलता येईल.}}%
}

%%K t13-tit-2 sub
\VUsaut{4.0mm}
\VUpk{10.2pt}{40}{उत्तम न्याहारी.}
\VUsaut{3.0mm}

%%K t13-09-1 p
९. --- वाढप्याने मला न्याहारीची पदार्थांची यादी (*) आणि वाइनची यादी आणून दिली. मी निवडले आणि
\VUgras{सोबतीचा पदार्थ} म्हणून \VUgras{लोण्यासोबत} \VUgras{कोळंबी} मागवली, आणि
\VUgras{पहिला पदार्थ} म्हणून \VUgras{काँच्या पद्धतीची आतडी}. मी \VUgras{मासा} खाल्ला नाही, पण
कोकराच्या \VUgras{मांडीचा} एक वाटा मागितला, आणि \VUgras{भाजीचा पदार्थ} म्हणून मलईसोबत
\VUgras{पालक}. नंतर \VUgras{मधल्या पदार्थांपैकी} एकही मला आवडला नाही, म्हणून मी वेगवेगळे गोड
पदार्थ आणवले : \VUgras{चीज}, \VUgras{फळे} आणि \VUgras{बिस्किटे}. शिवाय मी उत्तम सँ-एमिलिओं
\VUgras{वाइनही} घेतली, आणि जेवणाने फार समाधानी होऊन \VUgras{वाढप्याला} \textsuperscript{(57)}
चांगली \VUgras{बक्षिशी} \textsuperscript{(59)} देऊन मी मेज सोडले.

%%K t13-10-1 p
१०. --- मी निघायच्या तयारीत आहे हे पाहून \VUgras{व्यवस्थापकाने} \textsuperscript{(60)} मला
सुचवले की \VUgras{महासागराचा} \textsuperscript{(62)} भव्य देखावा पाहायला सज्ज्यावर जावे. दूरवर
मासेमारीच्या \VUgras{लहान होड्या} \textsuperscript{(63)}, \VUgras{शिडाची गलबते}
\textsuperscript{(64)} आणि एक अवाढव्य \VUgras{टपाल-नौका} \textsuperscript{(65)} दिसू लागली,
जिची काळी आकृती पांढऱ्या \VUgras{ढगांवर} \textsuperscript{(67)} उठून दिसत होती —
\VUgras{क्षितिजावरच्या} \textsuperscript{(66)}. हलक्या झुळुकीने \VUgras{पताका}
\textsuperscript{(70)} फडफडत होती, जी \VUgras{पाटीवर} \textsuperscript{(69)} रोवली आहे,
\VUgras{दालनाच्या} \textsuperscript{(68)}.

%%K t13-tit-3 sub
\VUsaut{3.0mm}
\VUpk{10.2pt}{40}{करमणुकी.}
\VUsaut{3.0mm}

%%K t13-11-1 p
११. --- तो देखावा इतका रोचक होता की मी ठरवले की उद्या \VUgras{द्विनेत्री दुर्बीण}
\textsuperscript{(71)} किंवा \VUgras{दुर्बीण} \textsuperscript{(72)} घेऊन कॅफेच्या गच्चीवर
चढायचे.

%%K t13-12-1 p
१२. --- सज्जा सोडून मी \VUgras{पेय} \textsuperscript{(73)} घ्यायला हॉटेलच्या कॅफेकडे गेलो,
\VUgras{बिलियर्डचा डाव} \textsuperscript{(75)} खेळत. न थांबता मी कॅफेचे दालन ओलांडले, जिथे
पुष्कळ गिऱ्हाइके \VUgras{बुद्धिबळ} \textsuperscript{(78)}, \VUgras{चेकर्स}
\textsuperscript{(79)}, \VUgras{डॉमिनोच्या फरशा} \textsuperscript{(80)}, \VUgras{बॅकगॅमन}
\textsuperscript{(81)} किंवा \VUgras{पत्ते} \textsuperscript{(82)} खेळत होती, आणि मी सरळ
\VUgras{बिलियर्डच्या खोलीत} \textsuperscript{(74)} चढलो.

%%K t13-13-1 p
१३. --- डाव संपल्यावर मी तळमजल्यावर उतरलो आणि पत्र लिहायला वाढप्याकडे \VUgras{लिफाफा}
\textsuperscript{(84)}, \VUgras{कागद} \textsuperscript{(83)}, \VUgras{टपाल तिकीट}
\textsuperscript{(85)}, \VUgras{लेखणी} \textsuperscript{(87)} आणि \VUgras{शाई}
\textsuperscript{(86)} मागितली.

%%K t13-13-2 p
मग मी एका \VUgras{गाडीवानाला} \textsuperscript{(92)} हाक मारली, ज्याची \VUgras{गाडी}
\textsuperscript{(88)} कॅफेसमोर थांबली होती. मी त्याला तासाचा किंवा एका फेरीचा दर विचारला, आणि
आमचे दरावर एकमत झाल्यावर त्याने मला \VUgras{गाडीचे दार} \textsuperscript{(89)} उघडून दिले आणि
आत बसवले. तो पुन्हा \VUgras{बैठकीवर} \textsuperscript{(91)} चढला, \VUgras{चाबकाने}
\textsuperscript{(93)} घोड्यांना भरभर हाकले, आणि आम्ही एका मोहक फेरफटक्याला निघालो.

\VUsaut{6.0mm}
\VUfilet{22mm}
